\documentclass[11pt]{article}
\usepackage[a4paper,margin=27mm]{geometry}
\usepackage{amsmath,amssymb,amsthm,booktabs,longtable,hyperref}
\newtheorem{theorem}{Theorem}
\newtheorem{proposition}[theorem]{Proposition}
\newtheorem{lemma}[theorem]{Lemma}
\newtheorem{corollary}[theorem]{Corollary}
\theoremstyle{definition}\newtheorem{definition}[theorem]{Definition}
\newcommand{\Aut}{\operatorname{Aut}}
\newcommand{\core}{\operatorname{core}}
\newcommand{\N}{\operatorname{N}}
\newcommand{\Stab}{\operatorname{Stab}}
\title{Faithful group actions on a bipartite graph}
\author{Private recovery of original L059}
\date{13 October 2026}
\begin{document}\maketitle
\begin{abstract}
A fixed connected bipartite graph carries a finite collection of faithful
rank two coset geometries.  Their classification is exactly a classification
of edge-transitive subgroups of the group preserving the two parts, up to
conjugacy in that group.  The correspondence retains the acting group and
the ordered stabilizers.  We give the correspondence and its action
invariants, and prove an exhaustive enumeration by descending through
maximal subgroups.  We apply it to the graph families considered by De Saedeleer, Leemans,
Mixer and Pisanski.  The retained calculations classify 4254 markings
on 1950 inputs, including the subdivisions added in the published
edition.  Every entry of the four printed tables is recovered, with
an explicit correction to the quartic input threshold and an explicit
choice of ordered parts.  The historical quartic identification uses
attributed classification premises; it is not a classification of all
semisymmetric quartic graphs.
\end{abstract}

\section{The objects and their kernels}
The source of this application is \cite{L059}.  Its structural results
belong to the classical relation between flag-transitive incidence
geometries and coset actions, which the source attributes to Tits and
subsequent work.  We retain that attribution.  The proofs of the kernel, marking correspondence, duality and local
invariant statements rework the project's existing phase177 argument;
they do not count as a second recovery.  The maximal-subgroup descent
and the subdivision extension lemma give the present organization
of the finite computation.  Both use elementary group-action facts;
we make no claim that those underlying facts are new.

All groups and graphs in the application are finite.  A marked coset
geometry is an ordered triple $(G;P,Q)$ of a group and two subgroups.
Its vertices are the left cosets of $P$ and of $Q$, kept as distinct
types; two vertices of different types are adjacent when their cosets
intersect.  A marked isomorphism is a group isomorphism that carries
the first displayed subgroup to the first, and the second to the second.
The action of $G$ is left multiplication on the cosets.

We first make explicit the elementary graph facts used in the
construction.  Suppose a group $G$ is edge-transitive on a connected
graph with at least one edge.  Fix an edge $\{u,v\}$.  Every vertex
is an endpoint of an edge, and hence belongs to the orbit of $u$ or
the orbit of $v$.  There are therefore at most two vertex orbits.
If there are two, each edge joins the two orbits and the graph is
bipartite.  If there is one, the group is vertex-transitive; it may
or may not be arc-transitive.  In the latter case the arcs form two
orbits exchanged by reversal.  Each has one arc per undirected edge.
If the common valency is $d$ and the number of vertices is $n$,
each orbit contains $nd/2$ arcs.  Vertex-transitivity makes the number
of arcs in that orbit with any prescribed initial vertex equal to
$d/2$.  Thus $d$ is even.  This proves the source's trichotomy and
even-valency statement.  The stronger last sentence of its
Proposition~1 is false for a specified group action: the cyclic
group of order three on a triangle is vertex- and edge-transitive
but has two orbits on arcs.

A connected bipartite graph has exactly two choices of ordered
bipartition.  Starting at a vertex, the parity of path length fixes
the part of every vertex, so every graph automorphism either
preserves both parts or exchanges them.  Its colour-preserving
subgroup is consequently the whole group or has index two.  These
are the connected assertions of the source's Propositions~4--6.
They also show the point- and line-transitivity assertion of
Proposition~7 in the present connected setting.  Indeed, if the full
group is edge-transitive, two vertices in the same part can be
mapped to one another by a full automorphism: in the two-orbit case
these parts are precisely the orbits; in the one-orbit case use
vertex-transitivity.  Such an automorphism preserves the parts, so
the colour-preserving subgroup is transitive on each part.

The disconnected assertion in Proposition~5 needs an additional
restriction.  For example take three disjoint copies of $K_{1,2}$,
colour the centres of two black and the third white.  Each component
has no colour-exchanging automorphism, since the centre has a
different degree from its leaves.  An automorphism can nevertheless
exchange a black-centred component with the white-centred component
while fixing the third component.  It then neither preserves both
global parts nor exchanges them.  Our use of the proposition is
confined to connected graphs.

\begin{proposition}\label{prop:kernel}
Set $B=P\cap Q$ and $N=\core_G(B)$.  The action on the union of the
two vertex types and the action on the edges both have kernel $N$.
The marking $(G/N;P/N,Q/N)$ has the same incidence graph and acts
faithfully on it.  The graph has $[G:\langle P,Q\rangle]$ connected
components.  In particular, it is connected precisely when
$\langle P,Q\rangle=G$.
\end{proposition}
\begin{proof}
An adjacent pair has the form $(gP,gQ)$: choose $g$ in the nonempty
intersection of its cosets.  Two choices $g,h$ give the same pair
precisely when $h^{-1}g\in P\cap Q$.  Thus the edge set is the
$G$-set $G/B$, and its kernel is $\core_G(B)$.
The kernels on the two types are $\core_G(P)$ and $\core_G(Q)$.
Their intersection is $\core_G(P\cap Q)$, since conjugation preserves
intersection.  Because $N$ lies in both stabilizers, the maps
$gP\mapsto(gN)(P/N)$ and $gQ\mapsto(gN)(Q/N)$ are bijections and
preserve nonempty intersections.  The quotient kernel is trivial.

Write $L=\langle P,Q\rangle$.  Moving from one edge to another
through their common vertex changes a representative by right
multiplication by an element of $P$ or $Q$.  Consequently every edge
in the component of $(P,Q)$ is represented by an element of $L$.
Conversely, a word in elements of $P$ and $Q$ gives a sequence of
edges sharing successive vertices, so every such representative is in
that component.  Translating by $g$ gives the component represented
by $gL$.  This proves the component formula.
\end{proof}

Here faithfulness is faithfulness on both types together.  Neither
individual vertex action is required to be faithful.  Proposition
\ref{prop:kernel} gives the source's Theorem 1, Corollary 1, and the
connectivity and flag-stabilizer formulas used in its construction.

\section{The classification over a fixed graph}
Let $X$ be a connected bipartite graph with ordered parts $X_0,X_1$,
each containing at least two vertices.  Put $A=\Aut^+(X)$, the group
of graph automorphisms preserving each part, and fix an edge
$e=\{b,w\}$ with $b\in X_0$, $w\in X_1$.  We consider subgroups
of $A$ that act transitively on the edges.

\begin{theorem}\label{thm:markings}
The assignment
\[
                  H\longmapsto(H;H_b,H_w)
\]
induces a bijection from the $A$-conjugacy classes of edge-transitive
subgroups of $A$ to the marked-isomorphism classes of faithful coset
geometries whose ordered incidence graph is isomorphic to $X$.
\end{theorem}
\begin{proof}
An edge-transitive subgroup $H$ is transitive on each vertex type,
because each vertex belongs to an edge.  The maps
$hH_b\mapsto h(b)$ and $hH_w\mapsto h(w)$ are bijections from
the two coset spaces to the vertex parts.  They send an intersecting
pair to an edge.  Conversely, every edge is an image of $e$ under
some $h\in H$, so every adjacent pair comes from intersecting cosets.
This is a graph isomorphism.  The action is faithful because $H$ is
an actual subgroup of the graph automorphism group.

Conversely, use a type-preserving graph isomorphism to transport the
faithful action of a marking to $X$.  Its image is a subgroup $H\le A$.
The coset action is transitive on adjacent pairs, by the first paragraph
of Proposition \ref{prop:kernel}, hence $H$ is edge-transitive.
Choosing the base edge gives exactly its two vertex stabilizers.

It remains to verify the equivalence relation.  Suppose
$H'=aHa^{-1}$ with $a\in A$.  Edge transitivity supplies $h'\in H'$
carrying $a(e)$ to $e$.  As both permutations preserve types, $h'a$
fixes $b$ and $w$ separately.  Conjugation by $h'a$ is therefore a
marked isomorphism.

For the converse, let $\phi:H\to H'$ carry $H_b$ to $H'_b$ and
$H_w$ to $H'_w$.  Define
\[
 f(h(b))=\phi(h)(b),\qquad f(h(w))=\phi(h)(w).
\]
The stabilizer conditions make both definitions well defined and
bijective.  The coset description just proved makes their union an
incidence-preserving bijection, with incidence-preserving inverse.
Thus $f\in A$.  Its definition implies $fhf^{-1}=\phi(h)$ on every
vertex, so $fHf^{-1}=H'$.  This proves both injectivity and surjectivity.
\end{proof}

The proof shows that conjugacy by the pointwise edge stabilizer $A_e$
gives the same equivalence classes on this collection of subgroups.
The faithful marking with acting group $A$ is the source's stable
marking.  A faithful marking is stable exactly when $H=A$: the
reverse construction from its graph produces $A$, and a proper
subgroup of a finite group cannot be isomorphic to it.

Without the faithful requirement there are infinitely many marked
realizations of the same graph.  From any realization form
$(G\times C_m;P\times C_m,Q\times C_m)$ for $m\ge1$.  The common
factor disappears from the coset graph, whereas the group orders
distinguish the markings.  This is the source's nonuniqueness result.

\section{Equivalence and duality}
An abstract group isomorphism and an isomorphism of the underlying
incidence graphs do not in general determine a marked isomorphism.
The sentence following Proposition 13 of \cite{L059} needs this
qualification.  For example, put
\[
 G=C_4\times C_2^2,\quad P=C_4\times1,\quad Q=1\times C_2^2.
\]
Both $(G;P,Q)$ and $(G;Q,P)$ have faithful incidence graph $K_{4,4}$:
the factors intersect trivially and every coset of one factor meets
every coset of the other.  No group automorphism exchanges $P$ and
$Q$, since only the former has an element of order four.  This is
the counterexample retained from phase177; its already existing Lean
file checks the stated factor obstruction, intersection and incidence
facts, and is not a formalization of the census algorithm.

Marked isomorphism is also weaker than simultaneous inner conjugacy.
In $C_2^2$, choose independent $a,b$.  An invertible linear map sends
$(\langle a\rangle,\langle b\rangle)$ to
$(\langle a+b\rangle,\langle b\rangle)$, while every inner
automorphism is trivial.  This proves the existence assertion of
the source's Proposition 14.  Its Mathieu-group example has the
same mechanism: an outer automorphism exchanges two classes of
subgroups of order 432 and preserves the relevant class of subgroups
of order 192.  Our explicit finite witness specifies the automorphism
and the conjugating element; an arbitrary conjugating element is
not needed for the assertion.

More explicitly, the retained permutation certificate in degree
twelve has $G=M_{12}$ of order $95040$, subgroups $P,Q$ of orders
$432,192$, and an automorphism $\phi$ with
$\phi(P)=P'$, $\phi(Q)=Q'=Q^h$.  The orbit lengths of $P$ and
$P'$ on the twelve points are respectively $(12)$ and $(3,9)$.
They therefore cannot be conjugate even in $S_{12}$, whereas
$\phi$ is a marked isomorphism.  Both pairs generate $G$, and
their intersection has trivial core.  The certificate also records
actual normal subgroups and complements giving
\[
 P\cong (C_3^2\rtimes Q_8)\rtimes S_3,
 \qquad Q\cong Q_8\rtimes S_4.
\]
The second subgroup orders and all intersections are checked on
these explicit permutations.  Thus the source-specific Mathieu
example is retained as well as the elementary existence argument.

Put $F=\Aut(X)$ and allow its elements to exchange the two parts.
\begin{proposition}\label{prop:duality}
For an edge-transitive $H\le A$, the type-exchanging elements of
$D=\N_F(H)_e$, where $e$ is fixed as a set, correspond to the
automorphisms of $H$ interchanging $H_b$ and $H_w$.  This correspondence
preserves orders.  In particular the least order of a marked duality
is the least permutation order among these elements of $D$.
\end{proposition}
\begin{proof}
Conjugation gives one direction.  Conversely, for an automorphism
$\phi$ exchanging the stabilizers, use
$h(b)\mapsto\phi(h)(w)$ and $h(w)\mapsto\phi(h)(b)$.
As in Theorem \ref{thm:markings}, these formulas give an automorphism
of $X$ normalizing $H$ and fixing $e$ as a set.

The conjugation homomorphism on $D$ is injective.  A type-preserving
element in its kernel fixes both endpoints of $e$ and commutes with
$H$, so it fixes every edge translate, hence every vertex.  If a
type-exchanging element were in its kernel, it would imply
$H_b=H_w$.  Connectedness gives $H=\langle H_b,H_w\rangle$, so
both vertex parts would be singletons, contrary to our assumptions.
Injectivity proves the assertion about orders.
\end{proof}

In the excluded one-edge graph, a permutation of the two types
must be retained as additional data: conjugation on the trivial
acting group does not detect it.  Every cubic or quartic input in
the source lies in the nondegenerate case above.

\section{Subdivisions and reuse of the group actions}
Let $X$ now be a connected cubic graph, and write $S(X)$ for its
subdivision: each original edge $\{u,v\}$ is replaced by the path
$u,e,v$ through a new vertex $e$.  Order the two parts by putting
original vertices first and edge vertices second.  Their degrees are
three and two respectively.

\begin{proposition}\label{prop:subdivision}
Restriction to the original vertices identifies
$\Aut(S(X))=\Aut^+(S(X))$ with $F=\Aut(X)$.  Under this identification,
edge transitivity on $S(X)$ is arc transitivity on $X$.  For an
arc-transitive subgroup $H\le F$ and an original edge $e=\{u,v\}$,
the corresponding marking has stabilizers
\[
 P=H_u,\qquad Q=H_{\{u,v\}},\qquad B=H_u\cap H_v.
\]
Consequently its invariant columns can be computed on $X$ itself.
\end{proposition}
\begin{proof}
Degrees force every automorphism of $S(X)$ to preserve its two
parts.  Two original vertices are adjacent in $X$ exactly when they
have a common edge vertex in $S(X)$.  Thus restriction gives an
automorphism of $X$.  Its action on each edge vertex is forced by
the two endpoints, and every automorphism of $X$ extends in this
way.  An edge of $S(X)$ is an incident pair $(u,\{u,v\})$, in
bijection with the arc $(u,v)$ of $X$.  This bijection is equivariant.
The stabilizer formulas follow, since fixing $u$ and the set
$\{u,v\}$ is equivalent to fixing $u$ and $v$ separately.
\end{proof}

For the four Max entries, the corresponding coset actions can
be computed directly as
\[
 P\curvearrowright X(u),\quad
 Q\curvearrowright\{u,v\},\quad
 H\curvearrowright V(X),\quad H\curvearrowright E(X).
\]
For example, $pB\mapsto p(v)$ identifies the first left coset
action, and $hQ\mapsto h(e)$ identifies the last.  The other two
maps use $u$.  These are equivariant bijections because the point
stabilizers are exactly $B,B,P,Q$, respectively.  They preserve
primitivity and double transitivity, including for actions with a
kernel.  The separate component \emph{Computing the subdivision
columns on the original graph} gives the complete argument in
the right-action convention of GAP and records the exact change
to the finite implementation.  No classified subgroup changes.

For a bipartite $X$ this computation can reuse its entire earlier
classification of edge-transitive subgroups.  Put $A=\Aut^+(X)$.
Since $F$ is arc-transitive in the symmetric cubic census,
$[F:A]=2$.

\begin{lemma}\label{lem:subdivision-reuse}
Every arc-transitive subgroup of $F$ has the form
\[
 H=\langle K,x\rangle,\quad K\le A\text{ edge-transitive},\quad
 x\in N_F(K)\setminus A,\quad x^2\in K.
\]
Conversely every such pair gives an arc-transitive $H$, and
$H\cap A=K$.  To list the $F$-conjugacy classes of $H$, it is enough
to use one $K$ from each previously listed $A$-conjugacy class,
test a transversal of $K$ in $N_F(K)$, and deduplicate the resulting
subgroups under $F$-conjugacy.
\end{lemma}
\begin{proof}
If $H$ is arc-transitive, it contains an element exchanging the two
parts.  Its subgroup $K=H\cap A$ is normal of index two.  Given two
edges, direct their arcs from the first part to the second.  An
element of $H$ carrying one directed arc to the other preserves
the parts, so lies in $K$.  Thus $K$ is edge-transitive.  Any
$x\in H\setminus A$ normalizes $K$, has square in $K$, and generates
$H$ together with $K$.

Conversely, the hypotheses make $K\cup Kx$ a subgroup, whose
intersection with $A$ is $K$.  The two orientations of edges are
separate transitive $K$-sets, and $x$ interchanges them.  Their union
is therefore transitive under $\langle K,x\rangle$.
Replacing $x$ by another element of its $K$-coset does not change
this subgroup or the square condition.  Finally, conjugating any
$K$ into its retained $A$-representative also conjugates the associated
$H$ and $x$.  The transversal construction therefore misses no
$F$-conjugacy class.  Deduplication gives precisely the equivalence
relation in Theorem~\ref{thm:markings}, applied to $S(X)$.
\end{proof}

\section{The invariant columns}
For a marking obtained from $H$, write $P=H_b$, $Q=H_w$ and
$B=P\cap Q$.  The four inclusions in the source's Max column are
\[
                       B<P,\quad B<Q,\quad P<H,\quad Q<H.
\]
For any inclusion $K<L$, act on $L/K$.  The entry is $2T$ if the
action is doubly transitive, $M$ if primitive but not doubly transitive,
and $X$ otherwise.  This makes sense even when the action has a kernel.

For completeness, primitivity is equivalent to maximality of $K$.
If $K<J<L$, then $J/K$ is a nontrivial block through the base coset.
Conversely, let $\Delta$ be such a block and let $J$ be its setwise
stabilizer.  Whenever $gK\in\Delta$, the blocks $g\Delta$ and
$\Delta$ meet at $gK$, hence agree.  Thus $g\in J$, and
$\Delta=J/K$.  The possible blocks are exactly the intermediate
subgroups.  Double transitivity is equivalently the condition that
$K$ has two orbits on $L/K$.

The four Norm entries test whether $H,P,Q,B$ respectively are
normal in $A$.  The MaxOrd entry is the greatest element order in
$H$, which can be computed using one representative from each
element-conjugacy class.  Proposition \ref{prop:duality} gives the
Dual entry.  Using arbitrary dualities of the unmarked graph would
lose the required normalizer condition.

\begin{corollary}
For a cubic graph the first two Max entries cannot be $X$.
For a quartic graph they cannot be $M$.
\end{corollary}
\begin{proof}
In the cubic case both inclusions have prime index three, so no
intermediate subgroup exists.  In the quartic case the faithful
local image is a transitive subgroup of $S_4$.  Its order is one of
$4,8,12,24$.  An order-four action is regular and has a block system
from any subgroup of order two.  An order-eight subgroup is a Sylow
two-subgroup and preserves a partition into two pairs.  The remaining
possibilities are $A_4$ and $S_4$, both doubly transitive.  Thus there
is no primitive local action that is not doubly transitive.
\end{proof}

\section{Exhaustive descent through maximal subgroups}
The classification need not enumerate every subgroup of $A$ before
testing edge transitivity.  The following elementary observation is
the relevant reusable part of the present computation.

\begin{theorem}\label{thm:descent}
Suppose $A$ is edge-transitive.  Start with the class of $A$.  For each retained subgroup $L$, take
representatives of all its conjugacy classes of maximal subgroups.
Retain those that are edge-transitive, except for classes already
retained up to conjugacy in $A$.  Continue until every retained
subgroup has been processed.  The output is exactly one representative
of every $A$-conjugacy class of edge-transitive subgroups.
\end{theorem}
\begin{proof}
All retained subgroups satisfy the required property, by construction.
Every proper step strictly decreases group order, and the finite
group $A$ has only finitely many subgroups, so the process terminates.

Let $K\le A$ be edge-transitive.  Choose a chain
\[
             K=L_s<L_{s-1}<\cdots<L_0=A
\]
in which each inclusion is maximal.  Every $L_i$ is edge-transitive,
because it contains $K$.  Inductively suppose the class of $L_i$
has a processed representative.  Conjugate the next subgroup along
with $L_i$ to place it in that representative.  A further conjugation
inside the representative puts it among the chosen maximal-subgroup
representatives.  It is therefore retained, or its $A$-class was
already retained.  In either case that class is eventually processed.
This proves by induction that the class of $K$ occurs.  The conjugacy
test prevents duplicate output classes.
\end{proof}

If $A$ itself is not edge-transitive, no subgroup can be edge-transitive,
and the required list is empty.  Otherwise an inexpensive necessary
condition is $|E(X)|\mid |L|$; it permits
discarding a maximal subgroup before computing its edge orbits.
Such a discard is exact by the orbit-stabilizer formula.  Failure of
edge transitivity is also inherited by every subgroup, which is why
no descendant of a discarded subgroup is needed.

Each retained node in our finite receipt records its actual permutation
generators and all maximal-subgroup representatives.  Each outgoing
step records either the failed order test, the failed edge-orbit test,
or the index of its retained $A$-class.  The last alternative is checked
by actual conjugacy in $A$, not merely by abstract isomorphism or group
order.  Exhaustiveness of the finite maximal-subgroup calculation and
of the full graph automorphism calculation is an explicit dependency
on GAP, GRAPE and nauty.  The direct nauty path also verifies every
reported generator against the edge set and checks the generated
group order in GAP; its completeness still uses nauty's algorithm.
This is a computational certificate with a written
completeness proof, not a claim of full Lean verification.

\section{Finite interfaces and historical scope}
The existing five-graph phase177 census has 69 markings: five over
$K_{3,3}$, 58 over $K_{4,4}$, and two each over the cube, the Heawood
graph, and the subdivision of the Petersen graph.  It retains the
source's corresponding exceptional examples and the nonstable
$A_5<S_5$ realization.  Its witness files are reused rather than
counted as new results.

The current primary catalogue gives a further explicit graph on 110
vertices, named $C4[110,7]$.  Its full group is of order 1320 and
preserves the two parts.  Exhaustive subgroup enumeration gives
exactly two markings, with Max entries $2TXXM$ and $2TXMM$.
These are the source's two named exceptional codes.  The finite
receipt records every edge and both stabilizer pairs, so the
displayed dihedral notation can be checked through actual group
orders instead of an unstated notation convention.

The acting groups and stabilizers are verified by actual group
isomorphisms.  Writing $D_{2m}$ for the dihedral group of order
$2m$, the two markings on 110 vertices are
\[
 (\operatorname{PGL}(2,11);S_4,D_{24}),\qquad
 (\operatorname{PSL}(2,11);A_4,D_{12}).
\]
Their Borel subgroups have orders six and three.  The source prints
$D_6$ and $D_3$; our statement uses the verified orders explicitly.
In particular both vertex parts have 55 elements.  The remaining
unique quartic Max code, $XX2T2T$, occurs on ten vertices with
$(C_5\rtimes C_4;C_4,C_4)$ and trivial Borel subgroup.

The corresponding named cubic cases are checked in the retained
five-graph data: the two dual $S_3\times C_3$ markings over
$K_{3,3}$, $(A_4;C_3,C_3)$ over the cube, and
$(C_7\rtimes C_3;C_3,C_3)$ over the Heawood graph.  The unique
$YYYY$ markings have elementary abelian acting groups $C_3^2$
and $C_2^4$ over $K_{3,3}$ and $K_{4,4}$.  The same certificates
check the unique cubic $YYNN$, $YNYN$ and $NNNY$ cases against
the ordered stabilizers.  This verifies the group descriptions
separately from the aggregate table counts.

Historical input is a separate obligation.  The archived 2011
Mini-Census homepage specifies at most 150 vertices.  The retained
2008 and 2014 tables both yield 418 bipartite inputs with vertex
stabilizer smaller than 100.  Thirty are half-arc-transitive and
cannot contribute a marking: if $F$ is vertex- and edge-transitive
but not arc-transitive, an edge stabilizer fixes its two endpoints;
it lies in the index-two group $A$.  Hence the $A$-orbit of the edge
has half the size of its $F$-orbit.  No subgroup of $A$ can be
edge-transitive.

The source also explicitly counts $K_{4,4}$, whose full vertex
stabilizer has order 144.  The literal cutoff therefore does not
include an example explicitly present in its table.  The 388 graphs
with actual vertex stabilizer below 100, together with $K_{4,4}$,
yield 1353 markings.  Eight additional graphs with vertex stabilizer
128 or 144 yield exactly 290 more.  Thus the reproduced table uses
397 graphs and 1643 markings, with the explicit scope
\[
 |X|\le150,\qquad |\Aut(X)_v|<200,
 \qquad X\text{ bipartite and edge-transitive under }\Aut^+(X).
\]
Since the two parts of a regular bipartite graph have the same size,
this stabilizer threshold is equivalently
$|\Aut^+(X)|/|X|<100$.  It is not the actual vertex-stabilizer
threshold stated in the paper.

Under this explicit scope all eleven Norm counts equal the printed
counts.  The twelve Max counts also agree after reversing the two
ordered parts of current graph $C4[128,55]$, named $SS[128,16]$.
It is semisymmetric: its two color orders need not define isomorphic
ordered graphs.  This reversal transfers three entries from $X2TXX$
to $2TXXX$ and changes no Norm count.  The stored witnesses retain
the original coloring, and the aggregation records this reversal
explicitly.  These calculations establish a reproducible numerical
scope; they do not establish why the original table used it.

Historical identification uses the separate, reviewed component
\emph{Identifying the historical quartic input set}.  Wilson and
Poto\v cnik state that their expanded census contains all previously
known edge-transitive tetravalent graphs up to 512 vertices
\cite[\S1]{Recipes}.  Acquiring that coverage assertion and the
historical census's distinctness and invariant assertions, each
historical set in the corrected scope is contained in the current
set.  The retained 2008 and 2014 tables each give 397 classes, while
direct checks of the current edge data give 397 pairwise
nonisomorphic classes in that scope.  Inclusion and equal finite
cardinality therefore identify the sets.  This argument does not
identify every old tag, order the two parts, or claim completeness
among all possible semisymmetric quartic graphs.  The 28 explicit
isomorphisms reconstructed from the archived BC, MSY and GPS2
definitions remain supplementary witnesses, not substitutes for
the stated coverage premise.

The cubic scope is distinguished from the date of the cited
classification.  The primary Foster site documents extensions from
768 to 2048 vertices and then to 10000.  We acquire the completeness
of its cubic symmetric and semisymmetric lists from the cited
classification, rather than attempting another graph enumeration.
The retained primary LCF records reconstruct all 796 symmetric
inputs through 2048 vertices and all 43 semisymmetric inputs through
768 vertices.  The two non-Hamiltonian symmetric graphs, the Petersen
and Coxeter graphs, have explicit separate constructions.  Direct
checks verify their simplicity, connectedness, valency and pairwise
nonisomorphism.  Of the 796 symmetric graphs, 714 are bipartite.
For 261 earlier bipartite inputs, explicit vertex permutations identify
the retained graph and its current LCF realization, so their completed
subgroup classifications are reused.

Every required graph now has a completed finite receipt.  The 714
bipartite symmetric graphs give 1082 faithful markings; all ten Max
and all six Norm entries agree with the printed cubic tables.
The 43 semisymmetric cubic graphs give 102 markings.  The final 2014
source explicitly includes subdivisions of symmetric cubic graphs,
an obligation absent from its arXiv v2 text.  All 796 such inputs
have been classified and give 1427 markings.  Proposition~\ref{prop:subdivision}
and Lemma~\ref{lem:subdivision-reuse} compute these on the original
vertex sets.  Every entry retains the acting permutation group,
its two stabilizers, their intersection, and the Max, Norm, Dual
and MaxOrd columns.  Two initial nonbipartite calculations reached
their time limits and made no mathematical conclusion.  Their
separate successful retries use the direct action calculation; both
initial receipts and the changed engines remain in the record.

\begin{table}[ht]
\centering
\begin{tabular}{lrr}
\toprule
Input family & Graphs & Faithful markings\\
\midrule
Bipartite symmetric cubic, $n\le2048$ & 714 & 1082\\
Semisymmetric cubic, $n\le768$ & 43 & 102\\
Subdivisions of symmetric cubic, original $n\le2048$ & 796 & 1427\\
Historical quartic, corrected scope & 397 & 1643\\
\midrule
Total & 1950 & 4254\\
\bottomrule
\end{tabular}
\caption{Completed classifications.  The scopes are stated above;
no large-stabilizer or unbounded classification is asserted.}
\end{table}

For the named historical quartic exceptions we do not need to
identify every old catalogue label.  Within each of the identified
historical and current sets, the pairs $(n,|\Aut(X)|)$ equal to
$(8,1152)$, $(10,240)$ and $(110,1320)$ each occur once.  Since
these are graph invariants, equality of the sets identifies these
three historical entries with the three explicit graphs used above.
This argument does not extend an invariant match to a graph
isomorphism without the already established inclusion and uniqueness.

\begin{table}[p]
\centering
\begin{minipage}[t]{.55\textwidth}
\centering
\begin{tabular}{rrrrr}
\toprule
$B<P$ & $B<Q$ & $P<H$ & $Q<H$ & Count\\
\midrule
2T & 2T & 2T & 2T & 4\\
2T & 2T & M & M & 21\\
2T & 2T & X & X & 315\\
2T & M & M & 2T & 1\\
2T & M & X & X & 32\\
M & 2T & 2T & M & 1\\
M & 2T & X & X & 32\\
M & M & 2T & 2T & 1\\
M & M & M & M & 88\\
M & M & X & X & 587\\
\bottomrule
\end{tabular}
\end{minipage}\hfill
\begin{minipage}[t]{.41\textwidth}
\centering
\begin{tabular}{rrrrr}
\toprule
$H$ & $P$ & $Q$ & $B$ & Count\\
\midrule
N & N & N & Y & 1\\
Y & N & N & N & 404\\
Y & N & N & Y & 674\\
Y & N & Y & N & 1\\
Y & Y & N & N & 1\\
Y & Y & Y & Y & 1\\
\bottomrule
\end{tabular}
\end{minipage}
\caption{Cubic symmetric Max and Norm counts.  Every displayed count equals the corresponding printed entry of \cite{L059}.  The cubic table uses the bipartite symmetric graphs through 2048 vertices.}
\end{table}
\begin{table}[p]
\centering
\begin{minipage}[t]{.55\textwidth}
\centering
\begin{tabular}{rrrrr}
\toprule
$B<P$ & $B<Q$ & $P<H$ & $Q<H$ & Count\\
\midrule
2T & 2T & 2T & 2T & 10\\
2T & 2T & M & M & 5\\
2T & 2T & X & X & 38\\
2T & X & M & M & 1\\
2T & X & X & 2T & 11\\
2T & X & X & M & 1\\
2T & X & X & X & 19\\
X & 2T & 2T & X & 11\\
X & 2T & X & X & 26\\
X & X & 2T & 2T & 1\\
X & X & M & M & 17\\
X & X & X & X & 1503\\
\bottomrule
\end{tabular}
\end{minipage}\hfill
\begin{minipage}[t]{.41\textwidth}
\centering
\begin{tabular}{rrrrr}
\toprule
$H$ & $P$ & $Q$ & $B$ & Count\\
\midrule
N & N & N & N & 365\\
N & N & N & Y & 268\\
N & N & Y & N & 4\\
N & N & Y & Y & 3\\
N & Y & N & N & 4\\
N & Y & N & Y & 3\\
Y & N & N & N & 466\\
Y & N & N & Y & 525\\
Y & N & Y & N & 2\\
Y & Y & N & N & 2\\
Y & Y & Y & Y & 1\\
\bottomrule
\end{tabular}
\end{minipage}
\caption{Quartic Max and Norm counts.  Every displayed count equals the corresponding printed entry of \cite{L059}.  The quartic threshold and ordered-part correction are explicit in the text.}
\end{table}


The published edition retains the problem about large stabilizers;
the earlier arXiv edition also asks about cyclic configurations.
A finite list and the classification theorem above settle neither
open-ended classification problem.

\begin{thebibliography}{9}
\bibitem{L059} J. De Saedeleer, D. Leemans, M. Mixer and T. Pisanski,
\emph{Core-free, rank two coset geometries from edge-transitive
bipartite graphs}, Mathematica Slovaca 64 (2014), 991--1006.
Fixed primary edition: arXiv:1106.5704v2, 29 June 2011.
\bibitem{C4} S. Wilson and P. Poto\v cnik,
\emph{A census of edge-transitive tetravalent graphs}, Mini-Census
and Full Census.  Exact archived editions and the current graph data
are recorded in the accompanying source ledger.
\bibitem{Recipes} S. Wilson and P. Poto\v cnik,
\emph{Recipes for Edge-Transitive Tetravalent Graphs},
arXiv:1608.04158v1 (2016).
\bibitem{Foster} M. Conder and P. Poto\v cnik,
\emph{Edge-transitive cubic graphs: analysis, cataloguing and
enumeration}, Journal of Algebra \textbf{685} (2025), 703--737;
accompanying \href{https://fostercensus.graphsym.net/}{census}.
\end{thebibliography}
\end{document}
